\section{The ordinary hook with a three-box tail}
\label{sec:hook-three-source}

\paragraph{Status, attribution and scope.}
New author derivation D-HK301--03. The complete author, root and
separate /root/imprimitive whole-proof and primary-source reviews have
passed, with distinct exact spectrum reconstructions.  The cyclic-average character-defect method
was supplied by the root's HD-AT301; the new hook character
bounds and constants are derived here.  The exact weighted-cycle
and actual quotient framework is the previously checked
D-ATT07--09.  No finite computation is treated as a universal
proof.  Throughout $n\ge15$, let $U$ be the standard
$(n-1)$-dimensional rational representation of $S_n$, and put
\[
 W=\Lambda^3U\simeq S^{(n-3,1,1,1)},\qquad
 a=\binom{n-1}{3},\quad
 A=\frac{n^2-9n+8}{6},\quad
 B=\frac{n^2-2n-1}{8},\quad
 b=\frac{(n-1)^2}{8},\quad d_0=A+B.
 \tag{HK3.1}
\]
These bounds are rational numbers; no integrality of $A,B,b$
is assumed.  For a permutation $x$, set
$I(x)=n-\#\operatorname{cycles}(x)$, including fixed letters.
For a complex operator $T$ on a space of dimension $a$, write
$\nu(T)=a-\max_z\dim\ker(T-zI)$.

\paragraph{Exact published representation and character inputs.}
Fulton--Harris, Proposition~3.12 and its complete proof,
pp.~31--32, give the irreducibility of every exterior power
of the standard representation
\cite[Proposition~3.12, pp.~31--32]{HK:FultonHarris1991}.
The hook identification is also described in Exercise~4.6,
p.~48, and the branching discussion on p.~59.
For an explicit published character formula with the precise
partition label, use Garsia--Goupil, Proposition~I.1, p.~2,
its complete proof (2.1)--(2.6), pp.~9--11, and Table~5.1,
p.~31 \cite[Proposition~I.1; Table~5.1]{HD:GarsiaGoupil2009}.
Writing $f,u,v$ for the numbers of one-, two-, and three-cycles,
their formula is
\[
 \chi_W(x)=\frac{(f-1)(f-2)(f-3)}6-(f-1)u+v.
 \tag{HK3.2}
\]
It agrees directly with the exterior-power trace identity
$\chi_{\Lambda^3U}(x)=
(\chi_U(x)^3-3\chi_U(x)\chi_U(x^2)+2\chi_U(x^3))/6$.
In particular the formula identifies the exterior cube and
the labeled irreducible character without an unproved labeling
assumption.  Its value at the identity is the degree in (HK3.1).
The label is not self-conjugate for $n\ge15$, so restriction
to $A_n$ is the ordinary irreducible module of the accepted
alternating-core interface.  Its rational $S_n$ extension
supplies all compatible outer automorphisms.

\subsection*{D-HK301. Uniform support with a full positive constant}

\begin{lemma}[D-HK301]
For every nonidentity $x\in S_n$,
\[
                  \nu(Wx)\ge A I(x)+B\ge d_0.
 \tag{HK3.3}
\]
This bounds every target eigenspace and is invariant under
arbitrary scalar multiplication of $Wx$.  For every $x$,
including the identity, $\chi_W(x)\ge-b$.
\end{lemma}

\begin{proof}
For $f=0$, (HK3.2) equals $-1+u+v\ge-1\ge-b$.
For $f\ge1$, its first term is nonnegative, including
$f=1,2,3$, when it vanishes.  Since $2u\le n-f$,
\[
 \chi_W(x)\ge-\frac{(f-1)(n-f)}2
             \ge-\frac{(n-1)^2}{8}=-b.
 \tag{HK3.4}
\]
This proves the global character lower bound.

Put
\[
 c=\frac{n^2-6n+8}{6},\qquad
 c_f=\frac{f^2+(n-6)f+n^2-6n+9}{6}.
 \tag{HK3.5}
\]
For $f\ge1$ the negative $u$ term may be discarded and
$v\le(n-f)/3$.  Thus
\[
 \chi_W(x)\le\frac{(f-1)(f-2)(f-3)}6+\frac{n-f}{3}.
 \tag{HK3.6}
\]
Subtracting this upper bound from $a$ and dividing by
$n-f$ gives exactly $c_f$ when $f<n$.
Moreover $6(c_f-c)=f^2+(n-6)f+1>0$.
For $f=0$, the constraint $2u+3v\le n$ gives
$\chi_W(x)=-1+u+v\le n/2-1$, and
$a-n/2+1=nc$.  The identity has zero defect and zero
moved-letter count.  Consequently
\[
                    a-\chi_W(x)\ge c(n-f)
 \tag{HK3.7}
\]
for every permutation, with the stronger coefficient $c_f$
available whenever $f\ge1$ and $x\ne1$.

Fix $x$ of order $L$, put $I=I(x)$, and let $m_z$ be
the multiplicity of a root $z$ in $Wx$.  The average number
of fixed letters of $x^j$ is the number of cycles of $x$.
Averaging (HK3.7) therefore gives $a-m_1\ge cI$.
Every character value is real and at least $-b$.
For $z^L=1$, $z\ne1$, the exact Fourier formula gives
\[
 \begin{split}
 m_z-m_1
 &=\frac1L\sum_{j=0}^{L-1}
       \chi_W(x^j)(\operatorname{Re}z^{-j}-1)\\
 &\le\frac bL\sum_{j=0}^{L-1}(1-\operatorname{Re}z^{-j})=b.
 \end{split}
 \tag{HK3.8}
\]
The same upper bound is trivial for $z=1$.
Hence $\nu(Wx)\ge cI-b$; this is an all-eigenvalue bound,
not an identification of the largest eigenspace with the fixed space.

If $I\ge(n-2)/2$, use $c-A=n/2$ to obtain
\[
 \nu(Wx)-AI\ge\frac n2I-b
     \ge\frac{n(n-2)}4-\frac{(n-1)^2}{8}=B.
 \tag{HK3.9}
\]
If $1\le I<(n-2)/2$, put $f_0=n-2I>2$.
Every power of $x$ fixes at least $f_0$ letters: a
nontrivial cycle of length $l$ moves at most $2(l-1)$
letters, and taking a power cannot decrease the fixed set.
The coefficient $c_f$ is increasing for $f\ge0$, because
$n\ge15$.  Applying (HK3.6)--(HK3.7) to each nonidentity
power improves the cyclic average to
$a-m_1\ge c_{f_0}I$.  Combining this with (HK3.8), the
desired bound is equivalent to
\[
 (n-f_0)\bigl(f_0^2+(n-6)f_0+3n+1\bigr)
                              \ge3n(n-2).
 \tag{HK3.10}
\]
After subtracting the right side, the polynomial is
\[
 p(f_0)=-f_0^3+6f_0^2+(n^2-9n-1)f_0+7n.
 \tag{HK3.11}
\]
It is concave for $2\le f_0\le n-2$.
Its values at the endpoints are
\[
                p(2)=2n^2-11n+14,\qquad
                p(n-2)=n^2-12n+34.
 \tag{HK3.12}
\]
Both are positive for all $n\ge15$, so the whole interval
has positive values.  This proves the small-index case,
and thus the complete lemma.
\end{proof}

\subsection*{D-HK302. One factor over every original quadratic field}

\begin{theorem}[D-HK302]
In the actual rank-nine sixfold inducing-source setup, a
single ordinary core $A_n$ acting as $\Lambda^3U$, $n\ge15$,
is impossible.  Every compatible scalar and outer action
is allowed, and the original imaginary quadratic field
is arbitrary.
\end{theorem}

\begin{proof}
The normalizer is scalar times the rational $S_n$ image.
That rational image meets the scalars trivially.  Indeed a
scalar image centralizes the faithful nontrivial $A_n$ image,
so its permutation centralizes $A_n$; in $S_n$ this forces
the identity.  The scalar and permutation coordinates
are consequently actual homomorphisms.  The latter gives
the genuine transitive $n$-letter quotient cover.

For the actual intermediate base of genus $h$, the Euler
identity is $6=2a(h-1)+D$.  Since $a\ge364$, genus at least
two is impossible.  Genus one has at most six genuine
marks and period-image rank at most six.  Rank nine
therefore forces $h=0$, $r\ge12$ and $D=2a+6$.
The generating product-one tuple in the actual transitive
$n$-letter quotient has total natural index at least
$2n-2$ by Riemann--Hurwitz.

For every genuine entry with nonidentity permutation,
its actual moving rank is at least $AI+B$, by D-HK301.
A genuine scalar has index zero and moving rank $a\ge B$,
so the same inequality holds.  Summing gives
\[
 \begin{split}
 D-2a&\ge2A(n-1)+12B-2a\\
     &=\frac{n^2-6n-13}{6}>6\qquad(n\ge15),
 \end{split}
 \tag{HK3.13}
\]
contradicting the actual Euler identity.
No conductor condition or inducing-field equality is used.
\end{proof}

\subsection*{D-HK303. Every transitive tensor count over non-Eisenstein fields}

\begin{theorem}[D-HK303]
Suppose the ordinary transitive normal core is $A_n^t$
on $W^{\otimes t}$, with $W=\Lambda^3U$, $n\ge15$,
and $t\ge2$.  If the character field of the actual
inducing representation contains an imaginary quadratic
field $K\ne\mathbf Q(\sqrt{-3})$, this representation
cannot give an actual period-rank-nine sixfold source.
All compatible scalar, outer and factor-permuting
actions are included.
\end{theorem}

\begin{proof}
Write $e=a^t$.  The rational tensor-wreath normalizer
has trivial scalar intersection by its faithful conjugation
on the simple core factors.  It gives actual scalar and
rational coordinates, and an actual transitive $nt$-letter
quotient.  Let the finite scalar image be $\mu_M$.
The inducing character field lies in $\mathbf Q(\zeta_M)$.
If that image were contained in $\mu_6$, it could not
contain the stated field $K$.  The generating product-one
inertia tuple thus contains at least two genuine marks
whose scalar coordinate lies outside $\mu_6$.
This does not assume equality of the inducing and original fields.

As before, the actual base has genus zero, $r\ge12$ and
$D=2e+6$.  A rational tensor-wreath operator has trace
equal to the product of the traces of its factor-cycle
products.  If this is negative, at least one factor has
absolute value at most $b$, and all others have absolute
value at most $a$.  The same applies to every power, so
all such traces are at least $-(b/a)e$.

At a scalar mark outside $\mu_6$, a fixed eigenvalue of
the actual operator would be a root of the rational
operator of order outside $1,2,3,6$.  If its order is
not four, its rational cyclotomic multiplicity is at
most $e/4$.  For order four the two primitive roots have
equal multiplicity $m$; the real trace of the square
is at most $e-4m$ and at least $-(b/a)e$.
Thus each of the two scalar marks moves at least
\[
                           \frac{3-b/a}{4}e.
 \tag{HK3.14}
\]
This proves the marked estimate directly with every scalar phase retained.

Every genuine entry moves at least $(d_0/a)e$ by
D-HK301 and convolution when factors are fixed.
A factor-moving entry has the proved weighted-word
support bound $(a-1)e/(2a)$.
For $n\ge15$ the elementary constants satisfy
\[
                  d_0\ge40,\qquad b\le d_0,\qquad
                  2d_0\le a-1,\qquad a\ge364.
 \tag{HK3.15}
\]
For example the numerators of $d_0-40$, $d_0-b$ and
$a-1-2d_0$, after substituting $n=x+15$, are respectively
$7x^2+168x+14$, $2x^2+42x+193$ and
$2x^3+71x^2+844x+3382$, with positive denominators.
Hence the lower $(d_0/a)e$ bound also covers moving factors.

If three entries move factors, choose the two scalar
marks and a third distinct factor-moving entry.
At least nine further genuine entries remain, and
\[
 D\ge2e+a^{t-1}\left(9d_0-\frac{b+1}{2}\right)>2e+6.
 \tag{HK3.16}
\]
This is impossible.  Transitivity rules out zero
factor-moving entries, and product one rules out one.
There are therefore exactly two.  Their factor
permutations are inverse, and transitivity makes each
a full $t$-cycle.

Let $H$ be the sum of the natural indices of the two
cycle products, and $N$ the sum of all within-factor
natural indices at the other $r-2\ge10$ entries.
The actual transitive $nt$-letter cover has total index
$2n(t-1)+H+N$, so Riemann--Hurwitz gives
\[
                             H+N\ge2n-2.
 \tag{HK3.17}
\]
At a factor-fixing entry with $s\ge1$ active factors,
the maximum of their projective supports is at least
their average.  D-HK301 therefore gives moving rank
at least $a^{t-1}(AI/s+B)\ge a^{t-1}(AI/t+B)$,
where $I$ is the entry's total natural index.
The entire positive constant $B$ is retained.
A genuine scalar satisfies the same formula with $I=0$.

For each weighted full $t$-cycle the complete general
lemma D-ATT08 gives
\[
 D_i\ge\left(1-\frac1t\right)e
       +\frac{a^{t-1}}t\nu(Wx_i)-\frac{R_t(a)}t.
 \tag{HK3.18}
\]
Only $\nu(Wx_i)\ge AI(x_i)$ is used here, so an identity
cycle product is included with both sides zero.
The exact error is $R_2(a)=a$, and for $t\ge3$ the
proved estimate is $R_t(a)\le2a^{t-2}$.
Consequently $2R_t(a)/a^{t-1}\le2$ for every $t\ge2$.
Summing the two cycles and the other entries gives
\[
 \begin{split}
 D-2e
 &\ge\frac{a^{t-1}}t
       \bigl(2A(n-1)-2a+10tB-2\bigr)\\
 &=\frac{a^{t-1}}t
       \left(10tB-\frac{2(n-1)(2n-1)}3-2\right)
       >B a^{t-1}>6.
 \end{split}
 \tag{HK3.19}
\]
For the strict step, $9tB\ge18B$ and
\[
 18B-\frac{2(n-1)(2n-1)}3-2
                   =\frac{11n^2-30n-59}{12}>0
 \tag{HK3.20}
\]
for all $n\ge15$.  This contradicts $D=2e+6$
and proves the theorem.
\end{proof}

\paragraph{Strictness and boundaries.}
This removes the entire hook-three supplier profile
in the stated ranges, not just a restatement of an
earlier numerical long-row condition.  For example
$(97,1,1,1)$ at $n=100$ has a long row beyond the
earlier Gaussian cutoff and is neither a standard nor
a two-box-tail label.  No actual source at that profile
was previously asserted.
The all-tensor result retains the non-Eisenstein
condition; the one-factor theorem does not need it.
No conclusion about a different label, a nontransitive
factor system or a positive algebraic-cycle supplier
is inferred.  These are actual finite-cover/representation
obstructions in the existing condition-A source gate.
No Markman, Severn or032 premise is introduced.
